% GENERATED FILE. Edit canonical lessons, then run npm run generate:books.
% canonical-source-hash: fnv1a64:41f493c8f7daa06b
% canonical-lessons: IT-C333-bollire, IT-C333-lessare, IT-C333-scolare, IT-C333-coperchio, IT-C333-tegame

\chapter{Sul fuoco}
\label{ch:it-sul-fuoco}
\hlchaptermodality{333}

\begin{chapteropening}
\textbf{By the end of this chapter:} \emph{I can say how I boiled and drained food, with the lid on the pan.}

Complete the last lesson of chapter 333: I can say how I boiled and drained food, with the lid on the pan.
\end{chapteropening}

\section[bollire]{bollire — to boil}

\label{lesson:IT-C333-bollire}



\begin{quote}
Before the new one: say the Italian for to peel, then the Italian for to grate.
\end{quote}

\subsection*{You'll want to know: bollire}
\textbf{bollire} — \textquotedblleft{}to boil\textquotedblright{}.

\begin{quote}
Quando l'acqua bolle, butta la pasta.
\end{quote}

\begin{quote}
\emph{When the water boils, put the pasta in.}
\end{quote}

Present \emph{bollo}, \emph{bolli}, \emph{bolle}.

\begin{cousinweb}
\emph{bollire} is inherited from Latin \emph{bullīre} \textquotedblleft{}to bubble, to boil\textquotedblright{} (EtymDB), from \emph{bulla} \textquotedblleft{}a bubble\textquotedblright{}.
\end{cousinweb}

\begin{grammarlens}[title={What you've built}]
One more word for telling, explaining and describing at B1.
\end{grammarlens}

\subsection*{Your turn}
\begin{itemize}
\raggedright
  \item \emph{Say it:} \emph{bollire}
  \item \emph{Say it:} \emph{bollire}, once more
  \item \emph{Say it:} a sentence of your own with \emph{bollire}
  \item \emph{Recall:} say the Italian for to peel, then the Italian for to grate, then say \emph{bollire} again
\end{itemize}

\begin{tcolorbox}[breakable,colback=teal!4,colframe=teal!35!black,title={Before you move on}]
What does bollire mean? (\textquotedblleft{}To boil\textquotedblright{}.) Say it once more.
\end{tcolorbox}
\section[lessare]{lessare — to cook in boiling water, to boil (food)}

\label{lesson:IT-C333-lessare}



\begin{quote}
Before the new one: say the Italian for to grate, then the Italian for to boil.
\end{quote}

\subsection*{You'll want to know: lessare}
\textbf{lessare} — \textquotedblleft{}to cook in boiling water, to boil (food)\textquotedblright{}.

\begin{quote}
Ho lessato le patate per venti minuti.
\end{quote}

\begin{quote}
\emph{I boiled the potatoes for twenty minutes.}
\end{quote}

\emph{La carne lessa} or \emph{il lesso} — boiled meat.

\begin{cousinweb}
\emph{lessare} is inherited from Late Latin \emph{ēlixāre} \textquotedblleft{}to boil\textquotedblright{} (EtymDB), built on \emph{lixa} \textquotedblleft{}hot water\textquotedblright{}.
\end{cousinweb}

\begin{grammarlens}[title={What you've built}]
One more word for telling, explaining and describing at B1.
\end{grammarlens}

\subsection*{Your turn}
\begin{itemize}
\raggedright
  \item \emph{Say it:} \emph{lessare}
  \item \emph{Say it:} \emph{lessare}, once more
  \item \emph{Say it:} a sentence of your own with \emph{lessare}
  \item \emph{Recall:} say the Italian for to grate, then the Italian for to boil, then say \emph{lessare} again
\end{itemize}

\begin{tcolorbox}[breakable,colback=teal!4,colframe=teal!35!black,title={Before you move on}]
What does lessare mean? (\textquotedblleft{}To cook in boiling water, to boil (food)\textquotedblright{}.) Say it once more.
\end{tcolorbox}
\section[scolare]{scolare — to drain}

\label{lesson:IT-C333-scolare}



\begin{quote}
Before the new one: say the Italian for to boil, then the Italian for to cook in boiling water.
\end{quote}

\subsection*{You'll want to know: scolare}
\textbf{scolare} — \textquotedblleft{}to drain\textquotedblright{}.

\begin{quote}
Scola la pasta al dente.
\end{quote}

\begin{quote}
\emph{Drain the pasta when it is al dente.}
\end{quote}

For pasta and vegetables: \emph{scolare la pasta}.

\begin{cousinweb}
\emph{scolare} is inherited from Vulgar Latin \emph{excōlāre} (EtymDB), built on Latin \emph{cōlāre} \textquotedblleft{}to filter\textquotedblright{}.
\end{cousinweb}

\begin{grammarlens}[title={What you've built}]
One more word for telling, explaining and describing at B1.
\end{grammarlens}

\subsection*{Your turn}
\begin{itemize}
\raggedright
  \item \emph{Say it:} \emph{scolare}
  \item \emph{Say it:} \emph{scolare}, once more
  \item \emph{Say it:} a sentence of your own with \emph{scolare}
  \item \emph{Recall:} say the Italian for to boil, then the Italian for to cook in boiling water, then say \emph{scolare} again
\end{itemize}

\begin{tcolorbox}[breakable,colback=teal!4,colframe=teal!35!black,title={Before you move on}]
What does scolare mean? (\textquotedblleft{}To drain\textquotedblright{}.) Say it once more.
\end{tcolorbox}
\section[il coperchio]{il coperchio — the lid}

\label{lesson:IT-C333-coperchio}



\begin{quote}
Before the new one: say the Italian for to cook in boiling water, then the Italian for to drain.
\end{quote}

\subsection*{You'll want to know: il coperchio}
\textbf{il coperchio} — \textquotedblleft{}the lid\textquotedblright{}.

\begin{quote}
Metti il coperchio sulla pentola.
\end{quote}

\begin{quote}
\emph{Put the lid on the pot.}
\end{quote}

Built on \emph{coprire}, ‘to cover’.

\begin{cousinweb}
\emph{coperchio} is inherited from Latin \emph{cooperculum} \textquotedblleft{}a lid\textquotedblright{} (EtymDB), from \emph{cooperīre} \textquotedblleft{}to cover\textquotedblright{}.
\end{cousinweb}

\begin{grammarlens}[title={What you've built}]
One more word for telling, explaining and describing at B1.
\end{grammarlens}

\subsection*{Your turn}
\begin{itemize}
\raggedright
  \item \emph{Say it:} \emph{il coperchio}
  \item \emph{Say it:} \emph{il coperchio}, once more
  \item \emph{Say it:} a sentence of your own with \emph{il coperchio}
  \item \emph{Recall:} say the Italian for to cook in boiling water, then the Italian for to drain, then say \emph{il coperchio} again
\end{itemize}

\begin{tcolorbox}[breakable,colback=teal!4,colframe=teal!35!black,title={Before you move on}]
What does il coperchio mean? (\textquotedblleft{}The lid\textquotedblright{}.) Say it once more.
\end{tcolorbox}
\section[il tegame]{il tegame — the shallow pan, the casserole}

\label{lesson:IT-C333-tegame}



\begin{quote}
Before the new one: say the Italian for to drain, then the Italian for the lid.
\end{quote}

\subsection*{You'll want to know: il tegame}
\textbf{il tegame} — \textquotedblleft{}the shallow pan, the casserole\textquotedblright{}.

\begin{quote}
Cuoci le uova nel tegame.
\end{quote}

\begin{quote}
\emph{Cook the eggs in the shallow pan.}
\end{quote}

A wide, low pan with two handles.

\begin{cousinweb}
\emph{tegame} is usually traced to Greek \emph{t\'{\={e}}ganon} \textquotedblleft{}a frying pan\textquotedblright{}, through a Latin form (EtymDB has no entry).
\end{cousinweb}

\begin{grammarlens}[title={What you've built}]
One more word for telling, explaining and describing at B1.
\end{grammarlens}

\subsection*{Your turn}
\begin{itemize}
\raggedright
  \item \emph{Say it:} \emph{il tegame}
  \item \emph{Say it:} \emph{il tegame}, once more
  \item \emph{Say it:} a sentence of your own with \emph{il tegame}
  \item \emph{Recall:} say the Italian for to drain, then the Italian for the lid, then say \emph{il tegame} again
\end{itemize}

\begin{tcolorbox}[breakable,colback=teal!4,colframe=teal!35!black,title={Before you move on}]
What does il tegame mean? (\textquotedblleft{}The shallow pan, the casserole\textquotedblright{}.) Say it once more.
\end{tcolorbox}
